\documentclass[10pt,a4paper]{amsart}
\usepackage[margin=19mm]{geometry}
\usepackage{amsmath,amssymb,mathtools}
\usepackage[scr=rsfs,cal=euler]{mathalfa}
\usepackage{xfrac}
\usepackage[unicode,hidelinks,bookmarks=false]{hyperref}
\newcommand{\ie}{i.e.\ }
\newcommand{\cf}{cf.\ }
\newcommand{\resp}{respectively\ }
\newcommand{\Subsection}{\subsection}
\providecommand{\nobreakdash}{\nobreak\hbox{-}}
\setlength{\emergencystretch}{2em}
\multlinegap0pt
\usepackage{amssymb}
\usepackage{bm}
\newtheorem{theorem}{Theorem}
\newtheorem{proposition}{Proposition}
\newcommand{\C}{\mathbb{C}}         % Complex numbers
\newcommand{\M}[1]{\bm{#1}}     % For matrices
\newcommand{\T}[1]{\mathcal{#1}}    % For tensors
\title{Hamiltonian and Symplectic Tensors in the T-product Algebra}
\author{Susana Lopez-Moreno, Taehyeong Kim}
\date{Source: arXiv:2605.20829v1 (CC BY 4.0)}
\begin{document}
\maketitle

\noindent\emph{Source and licence.} Hamiltonian and Symplectic Tensors in the T-product Algebra. Version arXiv:2605.20829v1 (CC BY 4.0), distributed by arXiv under the Creative Commons Attribution 4.0 International licence, http://creativecommons.org/licenses/by/4.0/, as recorded in the arXiv licence metadata for this exact version and shown on the abstract page https://arxiv.org/abs/2605.20829v1. Full licence notice: \texttt{licenses/arxiv-2605.20829-CC-BY-4.0.txt}.\par\smallskip

\noindent\textit{Editorial scope.} Counted unit: the article's proposition prop:T-Hamiltonian (source0.tex lines 277--281). The article's complete original proof of it is delivered with it (source0.tex lines 282--409, one \texttt{proof} environment of 128 lines). No mathematics is written, repaired or re-derived: the statement and the proof are the article's own lines, in the article's own notation, and no retained line is substituted or re-flowed.  Retained uncounted as the source-local context the proof needs: lines 232--235.  Nothing was omitted from any retained span, so the package carries no omission record.  No retained line contains a citation, so the wrapper carries no bibliography and no bibliography entry.  No retained line contains a figure environment, an image-inclusion command or drawing code, so no image is delivered and the package declares no asset dependency.  Editorial formatting only, in addition: an AMS \texttt{amsart} preamble that restates the article's own declarations and macros used by the retained lines, namely the environments \texttt{theorem}, \texttt{proposition} on separate counters, together with the standard packages those lines need.  The retained environments print in this document as Theorem 1, Proposition 1 in order of appearance, whereas the article prints the same items with the numbering of its own full document.  The complete native source of the article as distributed in the arXiv e-print for this exact version is bundled unchanged as \texttt{sources/201020/source0.tex}.
\begin{theorem}
\label{thm:hamiltonian_equiv}
A tensor $\T{H} \in \C^{2n \times 2n \times p}$ is T-Hamiltonian if and only if each of its frontal slices in the Fourier domain, $\widehat{\M{H}}^{(i)}$ for $i=1, \dots, p$, is a Hamiltonian matrix.
\end{theorem}

\begin{proposition}
\label{prop:T-Hamiltonian}
Let $\T{H} \in \C^{2n \times 2n \times p}$. 
$\T{H}$ is T-Hamiltonian if and only if each Fourier-domain slice $\widehat{\M{H}}^{(i)}$ can be written in the block form $\widehat{\M{H}}^{(i)} = \begin{bmatrix} \M{A}_i & \M{B}_i \\ \M{C}_i & -\M{A}_i^H \end{bmatrix}$, where $\M{A}_i, \M{B}_i, \M{C}_i \in \C^{n \times n}$ and the matrices $\M{B}_i$ and $\M{C}_i$ are Hermitian for each $i=1, \dots, p$.
\end{proposition}

\begin{proof}
Assume first that $\T{H}$ is T-Hamiltonian. By Theorem~\ref{thm:hamiltonian_equiv}, each matrix $\widehat{\M{H}}^{(i)}$ is Hamiltonian. Fix $i \in \{1,\dots,p\}$, and write
$$
\widehat{\M{H}}^{(i)} =
\begin{bmatrix}
\M{A}_i & \M{B}_i \\
\M{C}_i & \M{D}_i
\end{bmatrix},
$$
where $\M{A}_i,\M{B}_i,\M{C}_i,\M{D}_i \in \C^{n\times n}$. We compute
$$
\M{J}\widehat{\M{H}}^{(i)}
=
\begin{bmatrix}
\M{0} & \M{I}_n \\
-\M{I}_n & \M{0}
\end{bmatrix}
\begin{bmatrix}
\M{A}_i & \M{B}_i \\
\M{C}_i & \M{D}_i
\end{bmatrix}
=
\begin{bmatrix}
\M{C}_i & \M{D}_i \\
-\M{A}_i & -\M{B}_i
\end{bmatrix}.
$$
Since $\widehat{\M{H}}^{(i)}$ is Hamiltonian, the matrix $\M{J}\widehat{\M{H}}^{(i)}$ is Hermitian. Hence
$$
\left(
\begin{bmatrix}
\M{C}_i & \M{D}_i \\
-\M{A}_i & -\M{B}_i
\end{bmatrix}
\right)^H
=
\begin{bmatrix}
\M{C}_i & \M{D}_i \\
-\M{A}_i & -\M{B}_i
\end{bmatrix}.
$$
Computing the conjugate transpose on the left-hand side gives
$$
\begin{bmatrix}
\M{C}_i^H & -\M{A}_i^H \\
\M{D}_i^H & -\M{B}_i^H
\end{bmatrix}
=
\begin{bmatrix}
\M{C}_i & \M{D}_i \\
-\M{A}_i & -\M{B}_i
\end{bmatrix}.
$$
Equality of the corresponding blocks yields
$$
\M{C}_i^H = \M{C}_i, \qquad
-\M{A}_i^H = \M{D}_i, \qquad
\M{D}_i^H = -\M{A}_i, \qquad
\M{B}_i^H = \M{B}_i.
$$
The second identity gives $\M{D}_i = -\M{A}_i^H$, and the first and fourth identities show that $\M{C}_i$ and $\M{B}_i$ are Hermitian. Therefore
$$
\widehat{\M{H}}^{(i)} =
\begin{bmatrix}
\M{A}_i & \M{B}_i \\
\M{C}_i & -\M{A}_i^H
\end{bmatrix},
$$
with $\M{B}_i^H = \M{B}_i$ and $\M{C}_i^H = \M{C}_i$.

Conversely, assume that for each $i=1,\dots,p$ we have
$$
\widehat{\M{H}}^{(i)} =
\begin{bmatrix}
\M{A}_i & \M{B}_i \\
\M{C}_i & -\M{A}_i^H
\end{bmatrix},
$$
where $\M{B}_i^H = \M{B}_i$ and $\M{C}_i^H = \M{C}_i$. We compute
$$
\M{J}\widehat{\M{H}}^{(i)}
=
\begin{bmatrix}
\M{0} & \M{I}_n \\
-\M{I}_n & \M{0}
\end{bmatrix}
\begin{bmatrix}
\M{A}_i & \M{B}_i \\
\M{C}_i & -\M{A}_i^H
\end{bmatrix}
=
\begin{bmatrix}
\M{C}_i & -\M{A}_i^H \\
-\M{A}_i & -\M{B}_i
\end{bmatrix}.
$$
Taking the conjugate transpose gives
$$
\left(
\begin{bmatrix}
\M{C}_i & -\M{A}_i^H \\
-\M{A}_i & -\M{B}_i
\end{bmatrix}
\right)^H
=
\begin{bmatrix}
\M{C}_i^H & (-\M{A}_i)^H \\
(-\M{A}_i^H)^H & (-\M{B}_i)^H
\end{bmatrix}
=
\begin{bmatrix}
\M{C}_i^H & -\M{A}_i^H \\
-\M{A}_i & -\M{B}_i^H
\end{bmatrix}.
$$
Using $\M{C}_i^H = \M{C}_i$ and $\M{B}_i^H = \M{B}_i$, we obtain
$$
(\M{J}\widehat{\M{H}}^{(i)})^H
=
\begin{bmatrix}
\M{C}_i & -\M{A}_i^H \\
-\M{A}_i & -\M{B}_i
\end{bmatrix}
=
\M{J}\widehat{\M{H}}^{(i)}.
$$
Hence $\widehat{\M{H}}^{(i)}$ is Hamiltonian for every $i$. Theorem~\ref{thm:hamiltonian_equiv} now implies that $\T{H}$ is T-Hamiltonian.
\end{proof}

\end{document}
